% sec6_2.tex — Section 6.2: ARMA Processes

Having introduced the concept of the filter, we can study with ease and grace a class
of linear processes called ARMA processes, which are a parameterization of the
coefficients of $\text{MA}(\infty)$ processes.

\subsection*{AR(1) and Its MA($\infty$) Representation}

A \textbf{first-order autoregressive process} (\textbf{AR(1)}) satisfies the following stochastic
difference equation:
\begin{equation}
  y_t = c + \phi y_{t-1} + \varepsilon_t \quad \text{or} \quad
  y_t - \phi y_{t-1} = c + \varepsilon_t \quad \text{or} \quad
  (1 - \phi L) y_t = c + \varepsilon_t,
  \label{eq:6.2.1}
\end{equation}
where $\{\varepsilon_t\}$ is white noise. If $\phi \neq 1$, let $\mu \equiv c/(1 - \phi)$ and rewrite this equation as
\begin{equation}
  (y_t - \mu) - \phi \cdot (y_{t-1} - \mu) = \varepsilon_t \quad \text{or} \quad
  (1 - \phi L)(y_t - \mu) = \varepsilon_t.
  \label{eq:6.2.1prime}
  \tag{6.2.1$'$}
\end{equation}
As will be seen in a moment, $\mu$, not $c$, is the mean of $y_t$ if $y_t$ is covariance-stationary.
For this reason, we will call \eqref{eq:6.2.1prime} a deviation-from-the-mean form.
The moving average is on the successive values of $\{y_t\}$, not on $\{\varepsilon_t\}$. The difference
equation is called stochastic because of the presence of the random variable $\varepsilon_t$.

We seek a covariance-stationary solution $\{y_t\}$ to this stochastic difference equation.
The solution depends on whether $|\phi|$ is less than, equal to, or greater than~1.

\medskip
\noindent\textbf{Case 1:} $|\phi| < 1$

The solution can be obtained easily by the use of the inverse $(1 - \phi L)^{-1}$ given by
\eqref{eq:6.1.20}. Since this filter is absolutely summable when $|\phi| < 1$, we can apply it to
both sides of the AR(1) equation \eqref{eq:6.2.1prime} to obtain
\[
  (1 - \phi L)^{-1}(1 - \phi L)(y_t - \mu) = (1 - \phi L)^{-1} \varepsilon_t.
\]
This operation is legitimate because, thanks to Proposition~6.2, both sides of this
equation are well defined as mean square limits. By the definition \eqref{eq:6.1.14} of
inverses, $(1 - \phi L)(1 - \phi L)^{-1} = 1$, and by commutativity of inverses $(1 - \phi L)^{-1}(1 -
\phi L) = 1$. Therefore, if $\{y_t\}$ is covariance-stationary, the left-hand side equals $y_t - \mu$
by \eqref{eq:6.1.13}. So
\begin{equation}
  y_t - \mu = (1 - \phi L)^{-1} \varepsilon_t
  = (1 + \phi L + \phi^2 L^2 + \cdots) \varepsilon_t
  = \sum_{j=0}^{\infty} \phi^j \varepsilon_{t-j}
  \quad \text{or} \quad
  y_t = \mu + \sum_{j=0}^{\infty} \phi^j \varepsilon_{t-j}.
  \label{eq:6.2.2}
\end{equation}
What we have shown is that, if $\{y_t\}$ is a covariance-stationary solution to the stochastic
difference equation \eqref{eq:6.2.1} or \eqref{eq:6.2.1prime}, then $y_t$ has the \textbf{moving-average representation}
as in \eqref{eq:6.2.2}. Conversely, if $y_t$ has the representation \eqref{eq:6.2.2}, then
it satisfies the difference equation. This can be derived easily by substitution of
\eqref{eq:6.2.2} into \eqref{eq:6.2.1prime}:
\begin{align*}
  (1 - \phi L)(y_t - \mu) &= (1 - \phi L)(1 - \phi L)^{-1} \varepsilon_t \\
  &= \varepsilon_t \quad (\text{since } (1 - \phi L)(1 - \phi L)^{-1} = 1 \text{ by definition of inverses}).
\end{align*}
Thus, the process $\{y_t\}$ given by \eqref{eq:6.2.2} is the \emph{unique} covariance-stationary solution
to the first-order stochastic difference equation if $|\phi| < 1$. The mean of this process
is $\mu$ by Proposition~6.1(b).

The condition $|\phi| < 1$, which is the stability condition \eqref{eq:6.1.18} associated
with the first-order polynomial equation $1 - \phi z = 0$, is called the \textbf{stationarity
condition} in the present context of autoregressive processes. Intuitively, this result
says that the process, if it started a long time ago, ``settles down,'' provided that
$|\phi| < 1$ so that the effect of the past dies out geometrically as time progresses.

\medskip
\noindent\textbf{Case 2:} $|\phi| > 1$

By shifting time forward by one period (i.e., by replacing ``$t$'' by ``$t+1$''), multiplying
both sides by $\phi^{-1}$, and rearranging, the stochastic difference equation \eqref{eq:6.2.1prime}
can be written as
\begin{equation}
  y_t - \mu = \phi^{-1}(y_{t+1} - \mu) - \phi^{-1} \varepsilon_{t+1}.
  \label{eq:6.2.3}
\end{equation}
Applying the above argument but this time moving in the opposite direction in time
shows that the unique covariance-stationary solution is
\begin{equation}
  y_t = \mu - \sum_{j=1}^{\infty} \phi^{-j} \varepsilon_{t+j}.
  \label{eq:6.2.4}
\end{equation}
That is, the current value of $y$ is a moving average of \emph{future} values of $\varepsilon$. The infinite
sum is well defined because the sequence $\{\phi^{-j}\}$ is absolutely summable if $|\phi| > 1$.

\medskip
\noindent\textbf{Case 3:} $|\phi| = 1$

The stochastic difference equation has no covariance-stationary solution. For example,
if $\phi = 1$, the stochastic difference equation becomes
\begin{align*}
  y_t &= c + y_{t-1} + \varepsilon_t \\
  &= c + (c + y_{t-2} + \varepsilon_{t-1}) + \varepsilon_t \quad (\text{since } y_{t-1} = c + y_{t-2} + \varepsilon_{t-1}) \\
  &= c + \bigl(c + (c + y_{t-3} + \varepsilon_{t-2}) + \varepsilon_{t-1}\bigr) + \varepsilon_t, \quad \text{etc.}
\end{align*}
Repeating this type of successive substitution $j$ times, we obtain
\[
  y_t - y_{t-j} = c \cdot j + (\varepsilon_t + \varepsilon_{t-1} + \cdots + \varepsilon_{t-j+1}).
\]
(If $\{\varepsilon_t\}$ is independent white noise, this process is a \textbf{random walk with drift} $c$.)
Suppose, contrary to our claim, that the process is covariance-stationary. Then the
variance of $y_t - y_{t-j}$ is $2(\gamma_0 - \gamma_j)$. The variance of the right-hand side is $\sigma^2 \cdot j$. So
\[
  2 \cdot (\gamma_0 - \gamma_j) = \sigma^2 \cdot j \quad \text{or} \quad
  \rho_j = 1 - \frac{\sigma^2}{2\gamma_0} \cdot j < -1 \quad \text{for $j$ large enough}
\]
(recall: $\rho_j = \gamma_j / \gamma_0$).
This is a contradiction since autocorrelation coefficients cannot be greater than one
in absolute value. So $\{y_t\}$ cannot be covariance-stationary. Processes with $\phi = 1$
are examples of ``unit-root processes'' and will be studied in Chapter~9.

The solution \eqref{eq:6.2.4} for the case of $|\phi| > 1$, with the current value of $y$ linked
to the future values of the forcing process $\{\varepsilon_t\}$, is not useful in economics, which
does not confer perfect foresight on economic agents. In what follows, unless otherwise
stated, we will use the term ``an AR(1) process'' as the unique covariance-stationary
solution to an AR(1) equation \eqref{eq:6.2.1} or \eqref{eq:6.2.1prime} when the stationarity
condition holds.

\subsection*{Autocovariances of AR(1)}

The autocovariances $\{\gamma_j\}$ of AR(1) (for the case of $|\phi| < 1$) can be calculated in
two ways. The first method is to represent AR(1) as $\text{MA}(\infty)$ as in \eqref{eq:6.2.2} and use
\eqref{eq:6.1.7}. Since $\psi_j = \phi^j$ for AR(1), we have
\begin{equation}
  \gamma_j = (\phi^j + \phi^{j+2} + \phi^{j+4} + \cdots) \sigma^2
         = \frac{\phi^j \sigma^2}{1 - \phi^2}
  \quad \text{and} \quad
  \rho_j \;(\equiv \gamma_j / \gamma_0) = \phi^j
  \quad (j = 0, 1, \ldots).
  \label{eq:6.2.5}
\end{equation}
So the correlogram of an AR(1) process is very simple: it declines with $j$ at a
constant geometric rate.

The second method bases the calculation on what is called the \textbf{Yule--Walker
equations}; see Analytical Exercise~5 for the use of the Yule--Walker equations in
the calculation of autocovariances.

\subsection*{AR(\textit{p}) and Its MA($\infty$) Representation}

All the results just derived for AR(1) can be generalized to $\text{AR}(p)$, the \textbf{\textit{p}-th
order autoregressive process}, which satisfies the $p$-th order stochastic difference
equation:
\begin{equation}
\begin{gathered}
  y_t = c + \phi_1 y_{t-1} + \cdots + \phi_p y_{t-p} + \varepsilon_t \quad \text{or} \\
  y_t - \phi_1 y_{t-1} - \cdots - \phi_p y_{t-p} = c + \varepsilon_t \quad \text{or} \\
  \phi(L) y_t = c + \varepsilon_t \quad \text{with } \phi(L) = 1 - \phi_1 L - \phi_2 L^2 - \cdots - \phi_p L^p,
\end{gathered}
\label{eq:6.2.6}
\end{equation}
where $\{\varepsilon_t\}$ is white noise and $\phi_p \neq 0$. Evidently, $\phi(1) = 1 - \phi_1 - \cdots - \phi_p$. If
$\phi(1) \neq 0$, let $\mu \equiv c/(1 - \phi_1 - \cdots - \phi_p) = c/\phi(1)$. As shown below, $\mu$ is the
mean of $y_t$ if $y_t$ is covariance-stationary. Substituting for $c$, the AR($p$) equation
\eqref{eq:6.2.6} can be written equivalently in deviation-from-the-mean form:
\begin{equation}
  (y_t - \mu) - \phi_1 \cdot (y_{t-1} - \mu) - \cdots - \phi_p \cdot (y_{t-p} - \mu) = \varepsilon_t
  \quad \text{or} \quad \phi(L)(y_t - \mu) = \varepsilon_t.
  \label{eq:6.2.6prime}
  \tag{6.2.6$'$}
\end{equation}
The generalization to $\text{AR}(p)$ of what we have derived for AR(1) is

\begin{proposition}[AR(\textit{p}) as MA($\infty$) with absolutely summable coefficients]
\label{prop:6.4}
Suppose the $p$-th degree polynomial $\phi(z)$ satisfies the stationarity (stability) condition
\eqref{eq:6.1.18} or \eqref{eq:6.1.18prime}. Then
\begin{enumerate}[label=(\alph*)]
\item The unique covariance-stationary solution to the $p$-th order stochastic difference
equation \eqref{eq:6.2.6} or \eqref{eq:6.2.6prime} has the $\mathrm{MA}(\infty)$ representation
\begin{equation}
  y_t = \mu + \psi(L) \varepsilon_t, \quad
  \psi(L) = \psi_0 + \psi_1 L + \psi_2 L^2 + \psi_3 L^3 + \cdots\,,
  \label{eq:6.2.7}
\end{equation}
where $\psi(L) = \phi(L)^{-1}$. The coefficient sequence $\{\psi_j\}$ is bounded in absolute
value by a geometrically declining sequence and hence is absolutely summable.

\item The mean $\mu$ of the process is given by
\begin{equation}
  \mu = \phi(1)^{-1} c \quad \text{where $c$ is the constant in \eqref{eq:6.2.6}.}
  \label{eq:6.2.8}
\end{equation}

\item $\{\gamma_j\}$ is bounded in absolute value by a sequence that declines geometrically
with $j$. Hence, the autocovariances are absolutely summable.
\end{enumerate}
\end{proposition}

Unless otherwise stated, we will use the term ``an $\text{AR}(p)$ process'' as the unique
covariance-stationary solution \eqref{eq:6.2.7} to an $\text{AR}(p)$ equation \eqref{eq:6.2.6} that satisfies the
stationarity condition. The absolute summability of the inverse $\phi(L)^{-1}$ in part~(a)
of this proposition is immediate from Proposition~6.3. Given absolute summability,
the first part of~(a) can be proved by the same argument we used for AR(1); just
multiply both sides of \eqref{eq:6.2.6prime} by $\phi(L)^{-1}$ and observe that
\[
  \phi(L)^{-1} \phi(L) = \phi(L)\,\phi(L)^{-1} = 1.
\]
Part~(b) is immediate from Proposition~6.1(b). Part~(c) can easily be proved by
combining (a) and Proposition~6.1(b). Analytical Exercise~7 will ask you to give
an alternative proof of the proposition without using filters but using the apparatus
called the \textbf{companion form}.

As in AR(1), autocovariances can be obtained in two ways. The first method
utilizes the $\text{MA}(\infty)$ representation. The coefficients $\{\psi_j\}$ in the MA representation
are the coefficients in the inverse of the lag polynomial $\phi(L)$, and the algorithm for
calculating the inverse coefficients has already been presented (see \eqref{eq:6.1.16}). Given
$\{\psi_j\}$, use \eqref{eq:6.1.7} to calculate the autocovariances $\{\gamma_j\}$. The second method uses the
Yule--Walker equations.

\subsection*{ARMA(\textit{p}, \textit{q})}

An $\text{ARMA}(p, q)$ process combines $\text{AR}(p)$ and $\text{MA}(q)$:
\begin{equation}
\begin{gathered}
  y_t = c + \phi_1 y_{t-1} + \phi_2 y_{t-2} + \cdots + \phi_p y_{t-p} \\
  \quad + \theta_0 \varepsilon_t + \theta_1 \varepsilon_{t-1} + \theta_2 \varepsilon_{t-2} + \cdots + \theta_q \varepsilon_{t-q}
\end{gathered}
\end{equation}
or
\begin{equation}
\begin{gathered}
  y_t - \phi_1 y_{t-1} - \phi_2 y_{t-2} - \cdots - \phi_p y_{t-p} \\
  = c + \theta_0 \varepsilon_t + \theta_1 \varepsilon_{t-1} + \theta_2 \varepsilon_{t-2} + \cdots + \theta_q \varepsilon_{t-q}
\end{gathered}
\end{equation}
or
\begin{equation}
\begin{gathered}
  \phi(L) y_t = c + \theta(L) \varepsilon_t, \quad
  \phi(L) = 1 - \phi_1 L - \cdots - \phi_p L^p, \\
  \theta(L) = \theta_0 + \theta_1 L + \cdots + \theta_q L^q,
\end{gathered}
\label{eq:6.2.9}
\end{equation}
where $\{\varepsilon_t\}$ is white noise. If $\phi(1) \neq 1$, set $\mu = c/\phi(1)$. The deviation-from-the-mean
form is
\begin{equation}
\begin{gathered}
  (y_t - \mu) - \phi_1 \cdot (y_{t-1} - \mu) - \cdots - \phi_p \cdot (y_{t-p} - \mu) \\
  = \theta_0 \varepsilon_t + \theta_1 \varepsilon_{t-1} + \cdots + \theta_q \varepsilon_{t-q}
\end{gathered}
\end{equation}
or
\begin{equation}
  \phi(L)(y_t - \mu) = \theta(L) \varepsilon_t.
  \label{eq:6.2.9prime}
  \tag{6.2.9$'$}
\end{equation}
This is still a stochastic $p$-th order difference equation but with a serially correlated
forcing process $\{\theta(L) \varepsilon_t\}$ instead of a white noise process $\{\varepsilon_t\}$. Proposition~6.4
generalizes easily to

\begin{proposition}[ARMA(\textit{p},\,\textit{q}) as MA($\infty$) with absolutely summable coefficients]
\label{prop:6.5}
Suppose the $p$-th degree polynomial $\phi(z)$ satisfies the stationarity (stability)
condition \eqref{eq:6.1.18} or \eqref{eq:6.1.18prime}. Then
\begin{enumerate}[label=(\alph*)]
\item The unique covariance-stationary solution to the $p$-th order stochastic difference
equation \eqref{eq:6.2.9} or \eqref{eq:6.2.9prime} has the $\mathrm{MA}(\infty)$ representation
\begin{equation}
  y_t = \mu + \psi(L) \varepsilon_t, \quad
  \psi(L) = \psi_0 + \psi_1 L + \psi_2 L^2 + \psi_3 L^3 + \cdots\,,
  \label{eq:6.2.10}
\end{equation}
where $\psi(L) \equiv \phi(L)^{-1} \theta(L)$. The coefficient sequence $\{\psi_j\}$ is bounded in
absolute value by a geometrically declining sequence and hence is absolutely
summable.

\item The mean $\mu$ of the process is given by
\begin{equation}
  \mu = \phi(1)^{-1} c \quad \text{where $c$ is the constant in \eqref{eq:6.2.9}.}
  \label{eq:6.2.11}
\end{equation}

\item $\{\gamma_j\}$ is bounded in absolute value by a sequence that declines geometrically
with $j$. Hence, the autocovariances are absolutely summable.
\end{enumerate}
\end{proposition}

Unless otherwise stated, we will use the term ``an ARMA process'' as the unique
solution \eqref{eq:6.2.10} to an ARMA equation \eqref{eq:6.2.9} or \eqref{eq:6.2.9prime} that satisfies the
stationarity condition. The $\text{MA}(\infty)$ representation can be derived easily by multiplying
both sides of \eqref{eq:6.2.9prime} by $\phi(L)^{-1}$ and observing that $\phi(L)^{-1} \phi(L) =
\phi(L)\,\phi(L)^{-1} = 1$. For the rest of part~(a) and part~(b), the only nontrivial part
of the proof is to show that $\{\psi_j\}$ is bounded in absolute value by a geometrically
declining sequence. For $\text{AR}(p)$ we used the algorithm based on the set of equations
\eqref{eq:6.1.16} to solve $\phi(L)\,\psi(L) = 1$ for $\psi(L)$. This algorithm can easily be
generalized. This time $\psi(L) \equiv \phi(L)^{-1} \theta(L)$ or
\[
  \phi(L)\,\psi(L) = \theta(L).
\]
To solve this for $\psi(L)$, apply the convolution formula \eqref{eq:6.1.12} with $\alpha(L) = \phi(L)$,
$\beta(L) = \psi(L)$, and $\delta(L) = \theta(L)$ (instead of~1). The resulting set of equations is
the same as \eqref{eq:6.1.16}, except that the right-hand side of the equation for the constant
is $\theta_0$ rather than~1, that of the equation for $L$ is $\theta_1$ rather than~0, that of $L^2$ is $\theta_2$, and
so forth, until the equation for $L^q$, whose right-hand side is $\theta_q$. The right-hand side
of the rest of the equations is~0. As in $\text{AR}(p)$, these equations can easily be solved
successively for $(\psi_0, \psi_1, \psi_2, \ldots)$ as
\[
  \psi_0 = \theta_0 = 1, \quad \psi_1 = \theta_1 + \phi_1, \quad
  \psi_2 = \theta_2 + \phi_2 + \theta_1 \phi_1 + \phi_1^2, \quad \text{etc.}
\]
Again, as in the $\text{AR}(p)$ case, for sufficiently large $j$ (actually, for $j \geq \max(p,
q+1)$, as you will verify in Review Question~5), $\{\psi_j\}$ follows the same $p$-th order
homogeneous difference equation \eqref{eq:6.1.17}. So $\{\psi_j\}$ is bounded in absolute value
by a geometrically declining sequence and hence is absolutely summable.

Again there are two ways to derive the expression for autocovariances in terms
of $(\phi_1, \ldots, \phi_p, \theta_1, \ldots, \theta_q, \sigma^2)$. You can calculate the coefficients $\{\psi_j\}$ in the MA
representation as just described and then use \eqref{eq:6.1.7}. The other method is based on
the Yule--Walker equations.

\subsection*{ARMA(\textit{p}, \textit{q}) with Common Roots}

In the $\text{ARMA}(p, q)$ equation \eqref{eq:6.2.9prime}, suppose $\phi(L)$ satisfies the stationarity condition
as in Proposition~6.5, but suppose that one of the roots of $\phi(z) = 0$ is also a root
of $\theta(z) = 0$. If that root is $\lambda$, the polynomials $\phi(z)$ and $\theta(z)$ have a factorization
\[
  \phi(z) = \alpha(z)\,\phi^*(z), \quad
  \theta(z) = \alpha(z)\,\theta^*(z) \quad
  \text{with } \alpha(z) = 1 - \lambda^{-1} z.
\]
The inverse of $\phi(L)$ is $\phi^*(L)^{-1} \alpha(L)^{-1}$. So the $\psi(L)$ in Proposition~6.5 can be
written as
\begin{equation}
  \psi(L) = \phi(L)^{-1} \theta(L)
           = \phi^*(L)^{-1} \alpha(L)^{-1} \alpha(L)\,\theta^*(L)
           = \phi^*(L)^{-1} \theta^*(L),
  \label{eq:6.2.12}
\end{equation}
which shows that the $\text{ARMA}(p, q)$ equation \eqref{eq:6.2.9prime} and the simpler $\text{ARMA}(p-1,
q-1)$ equation
\[
  \phi^*(L)(y_t - \mu) = \theta^*(L) \varepsilon_t
\]
share the same process as the unique covariance-stationary solution.\footnote{Since we defined filters for sequences of real numbers, the discussion in the text presumes that the root $\lambda$ is real. If $\lambda$ is complex, then there must be another common root which is the complex conjugate of $\lambda$. Let $\lambda_1 = a + bi$ and $\lambda_2 = a - bi$ be the pair of complex roots. Then if we define
\[
  \alpha(z) = 1 - (\lambda_1^{-1} + \lambda_2^{-1})z + \frac{1}{(\lambda_1 \cdot \lambda_2)} z^2
  = 1 - \frac{2a}{(a^2 + b^2)} z + \frac{1}{(a^2 + b^2)} z^2,
\]
the discussion in the text carries over.}
For reasons of parsimony, ARMA equations with common roots are rarely used to parameterize
covariance-stationary processes.

\subsection*{Invertibility}

For the $\text{ARMA}(p, q)$ equation \eqref{eq:6.2.9}, if $\phi(z) = 0$ and $\theta(z) = 0$ have no common
roots and if $\theta(z)$ satisfies the stability condition \eqref{eq:6.1.18} or \eqref{eq:6.1.18prime}, then the
ARMA process is said to be \textbf{invertible} and the stability condition on $\theta(z)$ is called
the \textbf{invertibility condition}.\footnote{It is easy to confuse invertibility and the existence of an inverse. Since $\theta_0 = 1 \neq 0$, $\theta(L)$ does have an inverse. The inverse, however, may not be absolutely summable unless the invertibility condition is satisfied.}
Since $\theta(L)^{-1}$ is absolutely summable if $\theta(z)$ satisfies
the invertibility (stability) condition, we can multiply both sides of the ARMA
equation \eqref{eq:6.2.9} by $\theta(L)^{-1}$ and use the relation $\theta(L)^{-1} c = c\,\theta(1)^{-1} = c/\theta(1)$ to
obtain
\begin{equation}
  \theta(L)^{-1} \phi(L) y_t = \frac{c}{\theta(1)} + \varepsilon_t.
  \label{eq:6.2.13}
\end{equation}
So, the process has the infinite-order autoregressive ($\text{AR}(\infty)$) representation. The
derivation does not require the stationarity condition (that $\phi(z)$ satisfies the stability
condition). If both $\phi(z)$ and $\theta(z)$ satisfy the stability condition, then the
$\text{ARMA}(p, q)$ process has both the $\text{MA}(\infty)$ and $\text{AR}(\infty)$ representations.

\subsection*{Autocovariance-Generating Function and the Spectrum}

A particularly useful way to summarize the whole profile of autocovariances,
$\{\gamma_j\}$, of a covariance-stationary process $\{y_t\}$ is the \textbf{autocovariance-generating
function}:
\begin{equation}
  g_Y(z) \equiv \sum_{j=-\infty}^{\infty} \gamma_j z^j
  = \gamma_0 + \sum_{j=1}^{\infty} \gamma_j \cdot (z^j + z^{-j})
  \quad (\text{since } \gamma_{-j} = \gamma_j).
  \label{eq:6.2.14}
\end{equation}
The argument of this function ($z$) is a complex scalar. Because the summation
involves infinitely many terms, there arises a question as to whether the function
is well defined. A familiar result from calculus says that the infinite sum is well
defined at least for $|z| = 1$ (the unit circle) if $\{\gamma_j\}$ is absolutely summable. For
example, for $z = 1$,
\begin{equation}
  g_Y(1) = \sum_{j=-\infty}^{\infty} \gamma_j = \gamma_0 + 2 \sum_{j=1}^{\infty} \gamma_j.
  \label{eq:6.2.15}
\end{equation}
This is indeed well defined because the infinite sum converges.\footnote{Fact: If $\{\gamma_j\}$ is absolutely summable, then it is summable (i.e., the partial sum of $\{\gamma_j\}$ converges to a finite limit).}

Any complex number on the unit circle can be represented as
\begin{equation}
  z = \cos(\omega) - i \sin(\omega) = e^{-i\omega},
  \label{eq:6.2.16}
\end{equation}
where $i = \sqrt{-1}$ and $\omega$ is the negative of the radian angle that $z$ makes with the real
axis. If the autocovariance-generating function is evaluated at this $z$ and divided by
$2\pi$, the resulting function of $\omega$,
\begin{equation}
  s_Y(\omega) \equiv \frac{1}{2\pi}\,g_Y\bigl(\cos(\omega) - i\sin(\omega)\bigr)
  = \frac{1}{2\pi}\,g_Y(e^{-i\omega})
  \label{eq:6.2.17}
\end{equation}
is called the \textbf{(power) spectrum} or the \textbf{spectral density (function)} of $\{y_t\}$, and $\omega$
is referred to as the \textbf{frequency}.

It can be shown that for absolutely summable autocovariances, all of the autocovariances
can be calculated from the spectrum (hence the term autocovariance-generating
function). So there is a one-to-one mapping between the $\{\gamma_j\}$ sequence
and $g_Y(z)$ or $s_Y(\omega)$. The \textbf{time domain approach}, which concentrates on $\{\gamma_j\}$ and
which is what we have been doing, and the \textbf{frequency domain approach}, which
is based on the interpretation of the spectrum, are therefore \emph{equivalent}, although
some results can be more easily stated in one approach than in the other. This book
will not cover the frequency domain approach any further than this.

If $\{y_t\}$ is white noise, then evidently $g_Y(z)$ is constant and equal to the variance.
For MA(1), $g_Y(z)$ is obtained by simple substitution of \eqref{eq:6.1.2a} for $q = 1$ into the
definition \eqref{eq:6.2.14}:
\begin{align*}
  g_Y(z) &= \gamma_{-1} z^{-1} + \gamma_0 + \gamma_1 z \\
  &= (\theta_1 \sigma^2) z^{-1} + (1 + \theta_1^2)\sigma^2 + (\theta_1 \sigma^2) z \\
  &= \sigma^2 \cdot [\theta_1 z^{-1} + (1 + \theta_1^2) + \theta_1 z] \\
  &= \sigma^2 \cdot (1 + \theta_1 z)(1 + \theta_1 z^{-1}).
\end{align*}
So
\begin{equation}
  g_Y(z) = \sigma^2\,\theta(z)\,\theta(z^{-1}),
  \label{eq:6.2.18}
\end{equation}
where $\theta(z) = 1 + \theta_1 z$. This last expression generalizes to the $\text{MA}(q)$ case where
$\theta(z) = 1 + \theta_1 z + \theta_2 z^2 + \cdots + \theta_q z^q$. (You should verify this by carrying out the
multiplication in this expression for $g_Y(z)$ for $\text{MA}(q)$, collecting terms by powers
of $z$, and looking at the coefficient of $z^j$ or $z^{-j}$.) The expression is also good for
$\text{MA}(\infty)$ processes if the coefficients are absolutely summable, as the following
result assures us.

\begin{proposition}[Autocovariance-generating function for filtered processes]
\label{prop:6.6}
Let $\{\varepsilon_t\}$ be white noise and let $\psi(L) = \psi_0 + \psi_1 L + \psi_2 L^2 + \cdots$ be an absolutely
summable filter (i.e., with $\{\psi_j\}$ absolutely summable). Then the autocovariance-generating
function of the $\mathrm{MA}(\infty)$ process $\{y_t\}$ where $y_t = \mu + \psi(L) \varepsilon_t$ is
\begin{equation}
  g_Y(z) = \sigma^2\,\psi(z)\,\psi(z^{-1}).
  \label{eq:6.2.19}
\end{equation}
More generally, let $\{x_t\}$ be covariance-stationary with absolutely summable
autocovariances and $g_X(z)$ be the autocovariance-generating function of $\{x_t\}$. The
autocovariance-generating function of the filtered series $\{y_t\}$ where $y_t = h(L) x_t$
(with absolutely summable $\{h_j\}$) is
\begin{equation}
  g_Y(z) = h(z)\,g_X(z)\,h(z^{-1}).
  \label{eq:6.2.20}
\end{equation}
\end{proposition}

We will not prove this result; see, e.g., Fuller (1996, Theorem~4.3.1).

Because $\text{AR}(p)$ and $\text{ARMA}(p, q)$ processes have the $\text{MA}(\infty)$ representations,
their autocovariance-generating functions can be derived easily from \eqref{eq:6.2.19}.
Since $\psi(L) = 1/\phi(L)$ for $\text{AR}(p)$, and $\psi(L) = \theta(L)/\phi(L)$ for $\text{ARMA}(p, q)$, we
have
\begin{align}
  \text{AR}(p) &: \; g_Y(z) = \sigma^2 \frac{1}{\phi(z)\,\phi(z^{-1})},
  \label{eq:6.2.21} \\
  \text{ARMA}(p, q) &: \; g_Y(z) = \sigma^2 \frac{\theta(z)\,\theta(z^{-1})}{\phi(z)\,\phi(z^{-1})}.
  \label{eq:6.2.22}
\end{align}

\subsection*{Questions for Review}

\begin{enumerate}
\item If $\{y_t\}$ is the covariance-stationary solution to the AR(1) equation \eqref{eq:6.2.1} with
$|\phi| < 1$, what is
\[
  \widehat{\E}^*(y_t \mid 1, y_{t-1})
\]
(the least squares projection)? \textbf{Hint:} $\E(y_{t-1} \varepsilon_t) = 0$. Do the projection coefficients
depend on $t$? Is $\widehat{\E}^*(y_t \mid 1, y_{t-1}) = \E(y_t \mid y_{t-1})$? \textbf{Hint:} $\E(y_{t-1} \varepsilon_t) = 0$
but does it mean $\E(\varepsilon_t \mid y_{t-1}) = 0$? What is $\widehat{\E}^*(y_t \mid 1, y_{t-1}, y_{t-2})$? Now
suppose that $\{y_t\}$ is the covariance-stationary solution to the AR(1) equation
when $|\phi| > 1$. Is it true that $\E(y_{t-1} \varepsilon_t) = 0$? [Answer: No.] Does $\widehat{\E}^*(y_t \mid
1, y_{t-1}) = c + \phi y_{t-1}$ as was the case when $|\phi| < 1$?

\item (AR(1) with $\phi = -1$) \; For simplicity, set $c = 0$ in the AR(1) equation \eqref{eq:6.2.1}.
Show that the AR(1) equation with $\phi = -1$ has no covariance-stationary solution.
\textbf{Hint:} The same sort of successive substitution we did for the case $\phi = 1$
yields
\[
  y_t - (-1)^j y_{t-j} = \varepsilon_t - \varepsilon_{t-1} + \cdots + (-1)^{j-1} \varepsilon_{t-j+1}.
\]
From this, derive
\[
  (-1)^j \rho_j = 1 - \frac{\sigma^2}{2\gamma_0} \cdot j.
\]

\item (About Proposition~6.4) \; How do we know that $\phi(1) \neq 0$ in \eqref{eq:6.2.8}?
\textbf{Hint:} Use the stationarity condition. Prove part~(b) of Proposition~6.4. \textbf{Hint:} Take the expectation
of both sides of \eqref{eq:6.2.6} and exploit the covariance-stationarity of $\{y_t\}$.

\item (About Proposition~6.5) \; Verify that $y_t - \mu = \phi(L)^{-1}\,\theta(L) \varepsilon_t$ is a solution to
the $\text{ARMA}(p, q)$ equation.

\item ($\{\psi_j\}$ for $\text{ARMA}(p, q)$) \; For $\text{AR}(p)$, we used the equations \eqref{eq:6.1.12} to calculate
the MA coefficients $\{\psi_j\}$. For $\text{ARMA}(3, 1)$, write down the corresponding
equations. From which lag $j$ does $\{\psi_j\}$ start following
\[
  \psi_j - \phi_1 \psi_{j-1} - \phi_2\,\psi_{j-2} - \phi_3\,\psi_{j-3} = 0,
\]
which is \eqref{eq:6.1.17} for $p = 3$? Do the same for $\text{ARMA}(1, 2)$. From which $j$
does $\{\psi_j\}$ start following \eqref{eq:6.1.17} for $p = 1$? [Answer: $j = 3$.]

\item (How filtering alters the spectrum) \; Let $h(L)$ be absolutely summable. By
Proposition~6.2, if $\{x_t\}$ is covariance-stationary with absolutely summable autocovariances,
then so is $\{y_t\}$ where $y_t = h(L) x_t$. Using \eqref{eq:6.2.20}, verify that
\[
  s_Y(\omega) = h(e^{-i\omega})\,s_X(\omega)\,h(e^{i\omega}).
\]

\item (The spectrum is real valued) \; Show that the spectrum is real valued. \textbf{Hint:}
Substitute \eqref{eq:6.2.16} into \eqref{eq:6.2.14}. Use the following facts from complex analysis:
$[\cos(\omega) - i\sin(\omega)]^j = \cos(j\omega) - i\sin(j\omega)$; $\sin(-\omega) = -\sin(\omega)$.
\end{enumerate}
